\subsection{О сохранении фазового объёма.}
    Пусть $U_0 \subset \R^{2m}$ -- некоторая область в фазовом пространстве.
    Рассмотрим некоторую динамическую систему с гамильтонианом $H$.
    Она порождает фазовый поток, который для заданного $T$ диффеоморфно действует на $U_0$ и переводит в $U_T$.
    Справедлива следующая теорема:
    \begin{theorem}[Лиувилль, о сохранении фазового объёма]
        Для всякого $t$, для которого $U_t$ корректно определено, верно
        \[
            \mu(U_0) = \mu(U_t).
        \]

    \end{theorem}
    \begin{proof}
        Ранее показано, что фазовый поток является симплектическим диффеоморфизмом.
        Таким образом $\det D\phi^t = \pm 1$ и, следовательно, \[\mu(U_t) = \int_{U_t} d\bt{q}d\bt{p} = \int_{U_0} |\det D\phi| d\bt{q}d\bt{p} = \mu(U_0).\]
    \end{proof}
\begin{note}
    Здесь под <<$\mu$>> понимается мера Лебега подходящей размерности.
\end{note}
\subsection{Другие виды критериев каноничности преобразований и форм производящих функций.}
Пусть $\phi: (\bt{q}, \bt{p}) \mapsto (\tilde{\bt{q}}, \tilde{\bt{p}})$ -- некоторый диффеоморфизм.
\begin{theorem}
    Следующие утверждения эквивалентны:
    \begin{itemize}
        \item $\phi$ -- каноническое преобразование координат.
        \item $\forall i, k \hookrightarrow \{\tilde{q}_i, \tilde{q}_k\} = 0, \{\tilde{p_i}, \tilde{p_k}\} = 0, \{\tilde{q}_i, \tilde{p}_k\} = c\delta_{ik}$.
        \item $c\bt{p}\cdot\delta \bt{q} - \tilde{\bt{p}}\cdot \delta \tilde{\bt{q}} = \delta F$ (где под  <<$\delta$>> понимается изохронный полный дифференциал, а $F$ -- некоторая функция).
    \end{itemize}
\end{theorem}
\begin{proof}
    На самом деле, второе утверждение нам даёт вид <<матрицы Грама>> криволинейного скалярного произведения и он совпадает с $cJ$ -- чего нам необходимо и достаточно. \\
    Эквивалентность третьего утверждения получается тривиально если вспомнить о существовании такой вещи как производящая функция: $F$ ей и будет (в каких-либо подходящих координатах).
\end{proof}
На самом деле, вид производящей функции (и координаты, по которым мы её берём) можно варьировать в очень больших пределах.
Например, мы можем взять $S = S(\bt{q}, \tilde{\bt{q}}, t)$ и, если мы потребуем аналогичную невырожденность смешанного детерминанта, мы тоже придём к успеху.
Единственное отличие -- новые переменные будут генерироваться как
\[
    c\bt{p} = \dfrac{\partial S}{\partial \bt{q}}, \quad -\tilde{\bt{p}} = \dfrac{\partial S}{\partial \tilde{ \bt{q}}}.
\]
Также мы можем в производящей функции смешивать обобщённые координаты и импульсы (но, право дело, это абсолютно неудобно).
\subsection{Интегрируемость гамильтоновых систем.}
\begin{theorem}[Лиувилль, об интегрируемости в квадратурах]
    Пусть $2m$-мерная динамическая система с гамильтонианом $H$ имеет $m$ независимых первых интегралов в инволюции.
    Тогда она интегрируема в квадратурах.
\end{theorem}
\begin{proof}
    Пусть $f_i(\bt{q}, \bt{p}, t) = \alpha_i$ -- первые интегралы.
    Они попарно находятся в инволюции, то есть $\forall i, j \hookrightarrow\{f_i, f_j\} = 0$.
    И независимы, то есть матрица Якоби $\bt{f}$ имеет максимальный ранг.
    Поскольку матрица Якоби $\bt{f}$ имеет размер $m \times 2m$, то у неё существует невырожденный минор размера $m \times m$.
    Без ограничения общности будем считать, что этот минор является частью производной $\bt{f}$ по обобщённым импульсам ($\det D_{\bt{p}}\bt{f} \neq 0$) -- это влияет только на то, каким образом будет проинтегрирована система.
    Разрешим $\bt{p}$ относительно первых интегралов (что возможно в силу невырожденности $D_{\bt{p}}\bt{f}$).
    Тогда
    \[
        \bt{p} = \phi(\bt{q}, \bt{\alpha}, t).
    \]
    Подставим $\phi$ в $f_r$ и продифференцируем по $q_i$:
    \[
        \dfrac{\partial f_r}{\partial q_i} + \sum\limits_{k = 1}^m \dfrac{\partial f_r}{\partial p_k} \dfrac{\partial \phi_k}{\partial q_i} = 0.
    \]
    Умножим обе части на $\dfrac{\partial f_s}{\partial p_i}$ и просуммируем по $i$:
    \[
        \sum\limits_{i = 1}^m \dfrac{\partial f_r}{\partial q_i} \dfrac{\partial f_s}{\partial p_i} + \sum\limits_{i, k = 1, 1}^{m, m} \dfrac{\partial f_r}{\partial p_k} \dfrac{\partial f_s}{\partial p_i} \dfrac{\partial \phi_k}{\partial q_i} = 0.
    \]
    Поскольку полученное выражение верно для всяких $r$ и $s$, то оно верно и для $s, r$.
    Подставляя $s, r$ и вычитая одно из другого получим (с учётом инволютивности $f_r$ и $f_s$)
    \[
        \sum\limits_{i, k = 1, 1}^{m, m} \dfrac{\partial f_r}{\partial p_k}\dfrac{\partial f_s}{\partial p_i} \biggr(\dfrac{\partial \phi_k}{\partial q_i} - \dfrac{\partial \phi_i}{\partial q_k}\biggr) = 0.
    \]
    А значит $\dfrac{\partial \phi_i}{\partial q_k} = \dfrac{\partial \phi_k}{\partial q_i}$. \\
    Пусть $H^*$ -- функция Гамильтона в которой вместо $p_i$ подставлено $\phi_i$ и покажем, что
    \[
        \dfrac{\partial \phi_i}{\partial t} = - \dfrac{\partial H^*}{\partial q_i}.
    \]
    Поскольку
    \[
        -\dfrac{\partial H}{\partial q_i} = \dfrac{d p_i}{dt} = \dfrac{d\phi_i}{dt} = \dfrac{\partial \phi_i}{\partial t} + \sum\limits_{j = 1}^m \dfrac{\partial \phi_i}{\partial q_j}\dfrac{dq_j}{dt} = \dfrac{\partial \phi_i}{\partial t} + \sum\limits_{j = 1}^m \dfrac{\partial \phi_i}{\partial q_j} \dfrac{\partial H}{\partial p_j}.
    \]
    И тогда
    \[
        \dfrac{\partial \phi_i}{\partial t} = -\dfrac{\partial H}{\partial q_i} - \sum\limits_{j = 1}^m \dfrac{\partial \phi_j}{\partial q_i} \dfrac{\partial H}{\partial p_j} = - \dfrac{\partial H^*}{\partial q_i}.
    \]
    Таким образом существует такая $S(\bt{q}, \bt{\alpha}, t)$ что
    \[
        \dfrac{\partial S}{\partial q_i} = \phi_i, \quad \dfrac{\partial S}{\partial t} = -H^*.
    \]
    Она является решением уравнения
    \[
        dS = \sum\limits_{i = 1}^m \phi_i dq_i - Hdt.
    \]
    Её нахождение сведено к квадратурам.
    Функция $S$ является полным интегралом уравнения Гамильтона-Якоби и, следовательно, система проинтегрирована в квадратурах.
\end{proof}